\section{Introduction}
\IEEEPARstart{E}{nergy} benchmarking usually ends with a simple quantity: joules consumed by an execution, or the mean power needed to sustain it. Obtaining a trustworthy number requires several decisions that are less simple. A meter must observe the intended electrical boundary, cover the intended execution interval and represent the signal at an adequate rate. Repeating the experiment should also reproduce the work and operating conditions. A precise answer to only one of these questions can leave the final comparison misleading.

Consider an embedded power sensor, vendor telemetry and a plug-level AC meter. All three may produce stable readings during the same benchmark, yet they need not describe the same quantity. A Jetson module combines processing and supporting circuitry behind its DC input. A Hailo M.2 card performs inference with assistance from a separate host, which the card-input measurement excludes. A meter at the supply input adds conversion and any other connected loads. Agreement between these observations is therefore conditional on more than their nominal sampling rates.

The integration interval creates a second, independent difficulty. Two streams can report nearly the same average power while covering different parts of an execution. Their energies will then differ even before considering sensor calibration. Conversely, a low-rate trace can reproduce a long energy integral while missing most of the recorded power fluctuations. These are not contradictory results; they answer different measurement questions.

Our approach is to make those questions explicit and test them separately. High-rate recordings establish how a sampled representation changes an integral over fixed endpoints. Spectral analysis characterises the recorded fluctuations that an energy integral can average away. Physical repetitions then test the actual acquisition procedure, including interval selection, execution order and starting state. This sequence connects measurement validation to the settings used in a deployed benchmark rather than treating the oscilloscope result as a universal prescription.


The experimental study combines an integrated Jetson platform with a Hailo-10 M.2 accelerator. Jetson provides the controlled foundation: paired high-rate recordings quantify rate--duration decisions, and a separate long-window study tests whether energy agreement implies preserved fluctuations. Deployed PicoScope, \urecs{}, telemetry and scaled Shelly observations then test those decisions within each platform's recorded executions. The Hailo card-input study broadens the analysis to an attached accelerator with a different boundary and workload dynamics. We compare measurement procedures, not accelerator efficiency; the records do not establish identical compiled workloads across platforms.

Four research questions organise the complete measurement argument:
\begin{enumerate}
\item[\textbf{RQ1:}] Which tested sampling rates preserve interval energy, and how does the required rate change with the integration duration?
\item[\textbf{RQ2:}] How do energy-preserving rates differ from those needed to describe measured fluctuations, and how does the acquisition path affect that description?
\item[\textbf{RQ3:}] How does agreement between deployed measurement paths depend on the electrical boundary, workload and actual integration interval?
\item[\textbf{RQ4:}] How do execution duration and starting conditions affect the representativeness of repeated measurements after the rate has been chosen?
\end{enumerate}

The findings follow this sequence. Fixed-endpoint reconstruction identifies when low-rate energy integration is adequate. The same energy agreement can coexist with poor coverage of recorded fluctuations. Physical meter comparisons then expose a different issue: nearly equal mean powers can yield different energies when intervals or electrical domains differ. Finally, balanced acquisition and fixed-rate pause tests establish why execution conditions must be checked alongside the observation itself. Each step addresses a limitation that the preceding step cannot resolve alone.

This article extends our Jetson-focused study \emph{How Fast Is Fast Enough?}~\cite{mika2026workshopdraft}. Its measurement framework and central Jetson findings are presented here as the foundation of a self-contained journal account, not as newly obtained results. The additional experimental content comprises Hailo card-input and spectral measurements, duration-wise joint comparisons of energy, power and interval span, and a separate fixed-rate pause diagnostic. The instrumentation and initial validation also build on Wachsmuth's thesis~\cite{wachsmuth2026thesis}. The extension therefore broadens the measurement evidence while keeping its underlying quantities and validation logic explicit.
